% Options for packages loaded elsewhere
\PassOptionsToPackage{unicode}{hyperref}
\PassOptionsToPackage{hyphens}{url}
%
\documentclass[
]{article}
\usepackage{amsmath,amssymb}
\usepackage{iftex}
\ifPDFTeX
  \usepackage[T1]{fontenc}
  \usepackage[utf8]{inputenc}
  \usepackage{textcomp} % provide euro and other symbols
\else % if luatex or xetex
  \usepackage{unicode-math} % this also loads fontspec
  \defaultfontfeatures{Scale=MatchLowercase}
  \defaultfontfeatures[\rmfamily]{Ligatures=TeX,Scale=1}
\fi
\usepackage{lmodern}
\ifPDFTeX\else
  % xetex/luatex font selection
\fi
% Use upquote if available, for straight quotes in verbatim environments
\IfFileExists{upquote.sty}{\usepackage{upquote}}{}
\IfFileExists{microtype.sty}{% use microtype if available
  \usepackage[]{microtype}
  \UseMicrotypeSet[protrusion]{basicmath} % disable protrusion for tt fonts
}{}
\makeatletter
\@ifundefined{KOMAClassName}{% if non-KOMA class
  \IfFileExists{parskip.sty}{%
    \usepackage{parskip}
  }{% else
    \setlength{\parindent}{0pt}
    \setlength{\parskip}{6pt plus 2pt minus 1pt}}
}{% if KOMA class
  \KOMAoptions{parskip=half}}
\makeatother
\usepackage{xcolor}
\usepackage[margin=1in]{geometry}
\usepackage{color}
\usepackage{fancyvrb}
\newcommand{\VerbBar}{|}
\newcommand{\VERB}{\Verb[commandchars=\\\{\}]}
\DefineVerbatimEnvironment{Highlighting}{Verbatim}{commandchars=\\\{\}}
% Add ',fontsize=\small' for more characters per line
\usepackage{framed}
\definecolor{shadecolor}{RGB}{248,248,248}
\newenvironment{Shaded}{\begin{snugshade}}{\end{snugshade}}
\newcommand{\AlertTok}[1]{\textcolor[rgb]{0.94,0.16,0.16}{#1}}
\newcommand{\AnnotationTok}[1]{\textcolor[rgb]{0.56,0.35,0.01}{\textbf{\textit{#1}}}}
\newcommand{\AttributeTok}[1]{\textcolor[rgb]{0.13,0.29,0.53}{#1}}
\newcommand{\BaseNTok}[1]{\textcolor[rgb]{0.00,0.00,0.81}{#1}}
\newcommand{\BuiltInTok}[1]{#1}
\newcommand{\CharTok}[1]{\textcolor[rgb]{0.31,0.60,0.02}{#1}}
\newcommand{\CommentTok}[1]{\textcolor[rgb]{0.56,0.35,0.01}{\textit{#1}}}
\newcommand{\CommentVarTok}[1]{\textcolor[rgb]{0.56,0.35,0.01}{\textbf{\textit{#1}}}}
\newcommand{\ConstantTok}[1]{\textcolor[rgb]{0.56,0.35,0.01}{#1}}
\newcommand{\ControlFlowTok}[1]{\textcolor[rgb]{0.13,0.29,0.53}{\textbf{#1}}}
\newcommand{\DataTypeTok}[1]{\textcolor[rgb]{0.13,0.29,0.53}{#1}}
\newcommand{\DecValTok}[1]{\textcolor[rgb]{0.00,0.00,0.81}{#1}}
\newcommand{\DocumentationTok}[1]{\textcolor[rgb]{0.56,0.35,0.01}{\textbf{\textit{#1}}}}
\newcommand{\ErrorTok}[1]{\textcolor[rgb]{0.64,0.00,0.00}{\textbf{#1}}}
\newcommand{\ExtensionTok}[1]{#1}
\newcommand{\FloatTok}[1]{\textcolor[rgb]{0.00,0.00,0.81}{#1}}
\newcommand{\FunctionTok}[1]{\textcolor[rgb]{0.13,0.29,0.53}{\textbf{#1}}}
\newcommand{\ImportTok}[1]{#1}
\newcommand{\InformationTok}[1]{\textcolor[rgb]{0.56,0.35,0.01}{\textbf{\textit{#1}}}}
\newcommand{\KeywordTok}[1]{\textcolor[rgb]{0.13,0.29,0.53}{\textbf{#1}}}
\newcommand{\NormalTok}[1]{#1}
\newcommand{\OperatorTok}[1]{\textcolor[rgb]{0.81,0.36,0.00}{\textbf{#1}}}
\newcommand{\OtherTok}[1]{\textcolor[rgb]{0.56,0.35,0.01}{#1}}
\newcommand{\PreprocessorTok}[1]{\textcolor[rgb]{0.56,0.35,0.01}{\textit{#1}}}
\newcommand{\RegionMarkerTok}[1]{#1}
\newcommand{\SpecialCharTok}[1]{\textcolor[rgb]{0.81,0.36,0.00}{\textbf{#1}}}
\newcommand{\SpecialStringTok}[1]{\textcolor[rgb]{0.31,0.60,0.02}{#1}}
\newcommand{\StringTok}[1]{\textcolor[rgb]{0.31,0.60,0.02}{#1}}
\newcommand{\VariableTok}[1]{\textcolor[rgb]{0.00,0.00,0.00}{#1}}
\newcommand{\VerbatimStringTok}[1]{\textcolor[rgb]{0.31,0.60,0.02}{#1}}
\newcommand{\WarningTok}[1]{\textcolor[rgb]{0.56,0.35,0.01}{\textbf{\textit{#1}}}}
\usepackage{graphicx}
\makeatletter
\def\maxwidth{\ifdim\Gin@nat@width>\linewidth\linewidth\else\Gin@nat@width\fi}
\def\maxheight{\ifdim\Gin@nat@height>\textheight\textheight\else\Gin@nat@height\fi}
\makeatother
% Scale images if necessary, so that they will not overflow the page
% margins by default, and it is still possible to overwrite the defaults
% using explicit options in \includegraphics[width, height, ...]{}
\setkeys{Gin}{width=\maxwidth,height=\maxheight,keepaspectratio}
% Set default figure placement to htbp
\makeatletter
\def\fps@figure{htbp}
\makeatother
\setlength{\emergencystretch}{3em} % prevent overfull lines
\providecommand{\tightlist}{%
  \setlength{\itemsep}{0pt}\setlength{\parskip}{0pt}}
\setcounter{secnumdepth}{-\maxdimen} % remove section numbering
\ifLuaTeX
  \usepackage{selnolig}  % disable illegal ligatures
\fi
\usepackage{bookmark}
\IfFileExists{xurl.sty}{\usepackage{xurl}}{} % add URL line breaks if available
\urlstyle{same}
\hypersetup{
  pdftitle={TP 2 : Détection d'un changement ponctuel dans une suite binaire},
  pdfauthor={Hiba HANINI, Marwane HAJJY, Hiba FAHLI},
  hidelinks,
  pdfcreator={LaTeX via pandoc}}

\title{TP 2 : Détection d'un changement ponctuel dans une suite binaire}
\usepackage{etoolbox}
\makeatletter
\providecommand{\subtitle}[1]{% add subtitle to \maketitle
  \apptocmd{\@title}{\par {\large #1 \par}}{}{}
}
\makeatother
\subtitle{Modèles probabilistes pour l'apprentissage - MPA}
\author{Hiba HANINI, Marwane HAJJY, Hiba FAHLI}
\date{25/09/2025}

\begin{document}
\maketitle

\textbf{Q1. Soit} \(c = 1\). Donner l'expression de la loi générative
\(p(y \mid c = 1, \theta_1, \theta_2)\)

Comme les \(y_i\) sont i.i.d :

\[
\begin{aligned}
p(y \mid c = 1, \theta_1, \theta_2) &= \prod_{i=1}^n \theta_2^{y_i} (1 - \theta_2)^{1 - y_i}
\end{aligned}
\]

\textbf{Q2. Même question pour} \(c > 1\).

\[
\begin{aligned}
p(y \mid c = 1, \theta_1, \theta_2) &= \prod_{i=1}^{c-1} \theta_1^{y_i} (1 - \theta_1)^{1 - y_i} \prod_{i=c}^{n} \theta_2^{y_i} (1 - \theta_2)^{1 - y_i} 
\end{aligned}
\]

\textbf{Q3. On suppose que} \(\theta_1\) et \(\theta_2\) sont connues.
Déduire des questions précédentes que

\[
\forall c = 2, \ldots, n, \quad
\frac{p(c \mid y)}{p(c = 1 \mid y)} \;=\;
\prod_{j=1}^{c-1}
\frac{\theta_1^{y_j}(1 - \theta_1)^{1-y_j}}{\theta_2^{y_j}(1 - \theta_2)^{1-y_j}}.
\]

\textbf{Vérifier que l'on peut calculer le rapport}
\(p(c \mid y)/p(c = 1 \mid y)\) pour tout \(c\) en effectuant de l'ordre
de \(n(n - 1)/2\) multiplications.

On a d'après la formule de Bayes, \(\theta_1\) et \(\theta_2\) connues:

\[
\begin{aligned}
p(c \mid y) &= p(y \mid c) \cdot p(c)
\end{aligned}
\]

Et de même :

\[
\begin{aligned}
p(c = 1 \mid y) &= p(y \mid c = 1) \cdot p(c = 1)
\end{aligned}
\]

Comme \(p(c = 1) = p(c) = \frac{1}{n}\)

On obtient :

\[
\begin{aligned}
\forall c = 2, \ldots, n, \quad
\frac{p(c \mid y)}{p(c = 1 \mid y)} &= \frac{p(y \mid c)}{p(y \mid c = 1)} \\
&= \frac{\prod_{i=1}^{c-1} \theta_1^{y_i} (1 - \theta_1)^{1 - y_i} \prod_{i=c}^{n} \theta_2^{y_i} (1 - \theta_2)^{1 - y_i}}{\prod_{i=1}^n \theta_2^{y_i} (1 - \theta_2)^{1 - y_i}} \\
&= \prod_{j=1}^{c-1} \frac{\theta_1^{y_j}(1 - \theta_1)^{1-y_j}}{\theta_2^{y_j}(1 - \theta_2)^{1-y_j}}
\end{aligned}
\]

Si on calcule chaque rapport \(\frac{p(c \mid y)}{p(c = 1 \mid y)}\)
indépendamment en formant explicitement le produit pour chaque \(c\),
alors pour un \(c\) donné, il faut effectuer \((c - 1)\) multiplications
(en multipliant les \((c - 1)\) facteurs).

Le nombre total de multiplications pour tous les \(c = 2, \ldots, n\)
est donc :

\[
\sum_{c=2}^{n} (c - 1)
= 1 + 2 + \cdots + (n - 1)
= \frac{n(n - 1)}{2}.
\]

D'où la vérification, on peut bien calculer tous les rapports en
\(\mathcal{O}\big(\tfrac{n(n - 1)}{2}\big)\) multiplications.

\textbf{Q4. Supposant} \(\theta_1\) et \(\theta_2\) connues, proposer un
algorithme permettant de calculer \(p(c \mid y)\) pour tout
\(c = 1, \ldots, n\) avec une complexité en \(O(n)\).

n cherche à calculer \(p(c \mid y)\) pour tout \(c = 1, \ldots, n\) en
supposant que les paramètres \(\theta_1\) et \(\theta_2\) sont connus.

Soit : \[
r_j = 
\frac{\theta_1^{y_j}(1 - \theta_1)^{1 - y_j}}
{\theta_2^{y_j}(1 - \theta_2)^{1 - y_j}}.
\]

On peut alors écrire une récurrence linéaire pour mettre à jour le
rapport de manière itérative :

\[
R_1 = 1, \qquad
R_c = R_{c-1} \times r_{c-1} \quad \text{pour } c = 2, \ldots, n.
\]

Ainsi, on obtient directement : \[
\frac{p(c \mid y)}{p(c = 1 \mid y)} = R_c.
\]

Une fois tous les rapports \(R_c\) calculés, on peut normaliser pour
obtenir les probabilités \(p(c \mid y)\) :

\[
p(c \mid y) = \frac{R_c}{\sum_{k=1}^{n} R_k}.
\]

\begin{center}\rule{0.5\linewidth}{0.5pt}\end{center}

\subsubsection{✅ Algorithme en
pseudo-code}\label{algorithme-en-pseudo-code}

\begin{Shaded}
\begin{Highlighting}[]
\NormalTok{Entrée : y = (y1, ..., yn), θ1, θ2}
\NormalTok{Sortie : p(c | y) pour c = 1, ..., n}

\NormalTok{1. Initialiser R[1] ← 1}
\NormalTok{2. Pour j = 1 à n−1 :}
\NormalTok{       r\_j ← (θ1\^{}y[j] * (1−θ1)\^{}(1−y[j])) / (θ2\^{}y[j] * (1−θ2)\^{}(1−y[j]))}
\NormalTok{       R[j+1] ← R[j] * r\_j}
\NormalTok{3. Normaliser :}
\NormalTok{       somme ← Σ\_\{k=1\}\^{}\{n\} R[k]}
\NormalTok{       Pour c = 1 à n :}
\NormalTok{           p[c] ← R[c] / somme}
\NormalTok{4. Retourner p}
\end{Highlighting}
\end{Shaded}

\textbf{Q5. On suppose désormais que} \(\theta_1\) et \(\theta_2\) sont
inconnues. À partir de valeurs initiales arbitraires, proposer un
algorithme itératif, de type échantillonnage de Gibbs, permettant de
calculer \(p(c \mid y)\) pour tout \(c = 1, \ldots, n\) en combinant : -
la simulation de la loi \(p(c \mid y, \theta_1, \theta_2)\), - et la
simulation de valeurs des paramètres \(\theta_1\) et \(\theta_2\).

\textbf{Q6. Tester la convergence de votre algorithme en examinant la
sensibilité aux conditions initiales choisies arbitrairement.}

\textbf{Q7. Analyser les jeux de donn´ees envoy´es en pi`ece jointe par
l'enseignant. Décrire l'incertitude sur le(s) point(s) de changement
pour ces jeux de données (localisation des points des changements et
intervalles contenant chacun des points avec une probabilit´e
sup´erieure à 50\%).}

\end{document}
